\documentclass[11pt]{article}

\usepackage[margin=1in]{geometry}
\usepackage{amsmath, amssymb, amsthm}
\usepackage{booktabs}
\usepackage{hyperref}
\usepackage{microtype}

\newtheorem{theorem}{Theorem}
\newtheorem{corollary}[theorem]{Corollary}
\newtheorem{conjecture}[theorem]{Conjecture}
\newtheorem{definition}{Definition}
\newtheorem{remark}{Remark}

\newcommand{\Pm}{P_m}
\newcommand{\E}{\mathbb{E}}
\newcommand{\ord}{\mathrm{ord}}

\title{%
  Prime-Pair XOR Autocorrelation in the Binary Hypercube:\\
  A Hardy--Littlewood Analog
}

\author{Preetam Ojha \\ \texttt{nehu.preet@gmail.com}}
\date{May 2026}

\begin{document}
\maketitle

\begin{abstract}
Let $P_m$ denote the set of odd primes in $\{1,3,\ldots,2^m-1\}$.
For a parity-preserving mask $a$ (i.e.\ $a$ even, so $x \oplus a$ stays odd),
define the XOR autocorrelation
$C_m(a) = |\{(p,q) : p,q \in P_m,\; p \oplus q = a\}|$
and the ratio $\rho_m(a) = C_m(a)/\mathbb{E}_{\mathrm{rand}}[C_m(a)]$.
We define an \emph{XOR singular series}
$S_{\mathrm{xor}}(a) = \prod_{q \geq 3} \varphi_q(a)$
whose local factors $\varphi_q(a)$ measure the sieve obstruction at prime $q$.
We prove that for weight-1 masks $a = 2^j$, the limiting local factor is
$\varphi_q(2^j) = q(q-2)/(q-1)^2$, independent of~$j$,
and the infinite product equals the Hardy--Littlewood twin-prime constant
$C_2 \approx 0.66016$.
For weight-2 masks $a = 2^j + 2^k$, the local factor is
$q(2q-3)/(2(q-1)^2)$ when $2^{|k-j|} \equiv \pm 1 \pmod{q}$ (resonant case)
and $q(q-2)/(q-1)^2$ otherwise.
Numerical experiments at $m \in \{16, 18, 20\}$ confirm
$\rho_m(a) \approx S_{\mathrm{xor}}(a)$ with Pearson correlation $r \approx 0.88$
and mean ratio $\rho_m/S_{\mathrm{xor}} = 1.001 \pm 0.001$ across
thousands of masks of weight $1$--$4$.
We conjecture $\rho_m(a) \to S_{\mathrm{xor}}(a)$ as $m \to \infty$.
\end{abstract}

%% ---------------------------------------------------------------
\section{Setup}
%% ---------------------------------------------------------------

Fix $m \geq 2$ and let $Q_m^{\mathrm{odd}} = \{1, 3, 5, \ldots, 2^m - 1\}$
be the odd integers in $m$ bits, with $|Q_m^{\mathrm{odd}}| = 2^{m-1}$.
Write $P_m = \{p \in Q_m^{\mathrm{odd}} : p \text{ prime}\}$
and $\delta_m = |P_m|/2^{m-1}$.

\begin{definition}
A mask $a$ is \emph{parity-preserving} if $\mathrm{bit}_0(a) = 0$,
equivalently $a$ is even.
Then $x \oplus a$ is odd whenever $x$ is odd.
The \emph{XOR autocorrelation} at mask $a$ is the ordered-pair count:
\[
  C_m(a)
  \;=\;
  \sum_{x \in Q_m^{\mathrm{odd}}} \mathbf{1}[x \in P_m]\,\mathbf{1}[x \oplus a \in P_m].
\]
Each unordered prime pair $\{p,q\}$ with $p \oplus q = a$ contributes~2
to $C_m(a)$.
\end{definition}

Under a Bernoulli null model (primes are an independent $\delta_m$-subset),
$\mathbb{E}[C_m(a)] = \delta_m^2 \cdot 2^{m-1}$ for any $a \neq 0$.
Define the \emph{normalised ratio}:
\[
  \rho_m(a) \;=\; \frac{C_m(a)}{\delta_m^2 \cdot 2^{m-1}}.
\]

$C_m(a)$ is computable for all masks simultaneously via the
Walsh--Hadamard transform:
$C = \mathrm{IWHT}(\mathrm{WHT}(f)^2)$ where $f$ is the prime indicator
on $Q_m^{\mathrm{odd}}$, in time $O(m \cdot 2^m)$.

%% ---------------------------------------------------------------
\section{The XOR Singular Series}
%% ---------------------------------------------------------------

\begin{definition}
For an odd prime $q$ and a parity-preserving mask $a$, define the
\emph{local factor}:
\[
  \varphi_q(a)
  \;=\;
  \frac{P(q \nmid x \;\text{and}\; q \nmid x \oplus a)}{(1 - 1/q)^2},
\]
where the probability is over odd integers $x$ drawn uniformly
from $\{1, 3, \ldots, 2^m - 1\}$ as $m \to \infty$.
The \emph{XOR singular series} is:
\[
  S_{\mathrm{xor}}(a)
  \;=\;
  \prod_{\substack{q \;\text{prime} \\ q \geq 3}} \varphi_q(a).
\]
\end{definition}

\begin{remark}
For fixed bounded-popcount masks, $\varphi_q(a) = 1 + O(1/q^2)$
away from finitely many resonant primes (since $q(q-2)/(q-1)^2 = 1 - 1/(q-1)^2$),
so the Euler product converges absolutely.
\end{remark}

%% ---------------------------------------------------------------
\section{Main Results}
%% ---------------------------------------------------------------

\begin{theorem}[Weight-1 local factor]
\label{thm:weight1}
Let $j \geq 1$ and let $q$ be an odd prime.
In the limit $m \to \infty$, the local factor for mask $a = 2^j$ is:
\[
  \varphi_q(2^j)
  \;=\;
  \frac{q(q-2)}{(q-1)^2},
\]
independent of the bit position~$j$.
\end{theorem}

\begin{proof}
We need $P(q \nmid x \text{ and } q \nmid x \oplus 2^j)$ over odd $x$.
Write $x \oplus 2^j = x + 2^j(1 - 2\,\mathrm{bit}_j(x))$: if bit~$j$ of $x$ is 0,
the XOR adds $2^j$; if bit~$j$ is 1, it subtracts $2^j$.

For large $m$, the residue $x \bmod q$ and the bit $\mathrm{bit}_j(x)$
are asymptotically independent and jointly uniform, since $q$ is odd
and the binary representation imposes no congruence on bit~$j$.
Hence:
\[
  P(q \nmid x \oplus 2^j \mid x \equiv r \pmod{q})
  \;=\;
  \tfrac{1}{2} \cdot [r + 2^j \not\equiv 0] + \tfrac{1}{2} \cdot [r - 2^j \not\equiv 0]
  \pmod{q}.
\]
Since $2^j \not\equiv 0 \pmod{q}$ (as $q$ is odd), the two shifts
$r \mapsto r \pm 2^j$ hit $0 \bmod q$ at distinct values of $r$,
so at most one of the two conditions fails for any given~$r$.
Summing over $r \not\equiv 0 \pmod{q}$:
\[
  P(q \nmid x \text{ and } q \nmid x \oplus 2^j)
  \;=\;
  \frac{q-2}{q}.
\]
Dividing by the baseline $(1 - 1/q)^2 = (q-1)^2/q^2$ gives $q(q-2)/(q-1)^2$.
The result is independent of $j$ because the argument uses only
$2^j \not\equiv 0 \pmod{q}$, which holds for all $j \geq 1$
and all odd primes $q$.
\end{proof}

\begin{corollary}[Connection to the twin-prime constant]
\label{cor:C2}
For any bit position $j \geq 1$:
\[
  S_{\mathrm{xor}}(2^j)
  \;=\;
  \prod_{q \geq 3} \frac{q(q-2)}{(q-1)^2}
  \;=\;
  C_2
  \;\approx\;
  0.66016,
\]
where $C_2$ is the Hardy--Littlewood twin-prime constant.
\end{corollary}

\begin{proof}
The product identity $\prod_{q \geq 3} q(q-2)/(q-1)^2 = C_2$
is the standard Euler product representation of $C_2$
\cite{hardylittlewood1923}.
\end{proof}

\begin{theorem}[Weight-2 local factor]
\label{thm:weight2}
Let $j < k$, both $\geq 1$, set $d = k - j$, and let $q$ be an odd prime.
In the limit $m \to \infty$, the local factor for $a = 2^j + 2^k$ is:
\[
  \varphi_q(2^j + 2^k)
  \;=\;
  \begin{cases}
    \dfrac{q(2q-3)}{2(q-1)^2}
    & \text{if } 2^d \equiv \pm 1 \pmod{q}, \\[8pt]
    \dfrac{q(q-2)}{(q-1)^2}
    & \text{otherwise.}
  \end{cases}
\]
\end{theorem}

\begin{proof}
Flipping bits $j$ and $k$ maps $x$ to $x \oplus 2^j \oplus 2^k$.
Decompose on whether bits $j$ and $k$ of $x$ are equal or unequal
(each case has probability $1/2$ asymptotically, independently of $x \bmod q$):

\smallskip
\noindent\textbf{Case 1: resonant ($2^d \equiv \pm 1 \pmod{q}$).}
Here $2^k \equiv \varepsilon 2^j \pmod{q}$ for $\varepsilon \in \{+1,-1\}$.
\begin{itemize}
\item Bits $j = k$ in $x$ (probability $1/2$): the two flips contribute
  $\pm 2^j \pm \varepsilon 2^j = 0$ or $\pm 2^{j+1} \pmod{q}$.
  Specifically, when $\varepsilon = +1$ the shifts cancel ($x \oplus a \equiv x$);
  when $\varepsilon = -1$ they also cancel by similar symmetry.
  In either sub-case, $x \oplus a \equiv x \pmod{q}$,
  so $P(q \nmid x \oplus a \mid \text{bits } j = k) = P(q \nmid x) = (q-1)/q$.
\item Bits $j \neq k$ in $x$ (probability $1/2$): the net shift is
  $\pm 2^{j+1}$ or $\pm 2(2^j) \not\equiv 0$, so
  $P(q \nmid x \oplus a \mid \text{bits } j \neq k,\, q \nmid x) = (q-2)/(q-1)$.
\end{itemize}
Combining:
\[
  P(q \nmid x \text{ and } q \nmid x \oplus a)
  = \tfrac{1}{2} \cdot \tfrac{q-1}{q} + \tfrac{1}{2} \cdot \tfrac{q-2}{q}
  = \frac{2q-3}{2q}.
\]
Dividing by $(q-1)^2/q^2$ gives $q(2q-3)/(2(q-1)^2)$.

\smallskip
\noindent\textbf{Case 2: non-resonant ($2^d \not\equiv \pm 1 \pmod{q}$).}
The two shifts $\pm 2^j$ and $\pm 2^k$ are distinct and non-cancelling mod $q$,
so for each assignment of bits $j$, $k$ the combined shift is non-zero.
A residue computation analogous to Theorem~\ref{thm:weight1}
gives $P(q \nmid x \text{ and } q \nmid x \oplus a) = (q-2)/q$,
yielding the weight-1 factor $q(q-2)/(q-1)^2$.
\end{proof}

\begin{remark}[Resonance and multiplicative order]
The condition $2^d \equiv \pm 1 \pmod{q}$ holds if and only if
$\ord_q(2)$ divides $d$ or $\ord_q(2)$ divides $2d$ (but not $d$).
For $q = 3$, $\ord_3(2) = 2$, so $2^d \equiv \pm 1 \pmod 3$ for
\emph{all} $d \geq 1$: every weight-2 mask is resonant at $q = 3$,
with $\varphi_3 = 9/8$ universally.
\end{remark}

\begin{corollary}[Optimal weight-2 mask]
A weight-2 mask $2^j + 2^k$ maximises $S_{\mathrm{xor}}$ when
the bit-difference $d = k - j$ simultaneously satisfies
$2^d \equiv \pm 1 \pmod{q}$ for the largest possible set of small primes $q$.
This is equivalent to maximising the number of small prime factors of
$2^d - 1$ and $2^d + 1$.
For $d = 12$: $2^{12} - 1 = 3^2 \cdot 5 \cdot 7 \cdot 13$
and $2^{12} + 1 = 17 \cdot 241$,
giving simultaneous resonance at $q \in \{3, 5, 7, 13, 17\}$
and the largest predicted $S_{\mathrm{xor}}$ among $d \leq 20$.
\end{corollary}

%% ---------------------------------------------------------------
\section{The Conjecture}
%% ---------------------------------------------------------------

\begin{conjecture}[XOR Hardy--Littlewood]
\label{conj:main}
For any fixed parity-preserving mask $a$ with $\mathrm{popcount}(a)$ bounded:
\[
  C_m(a)
  \;\sim\;
  S_{\mathrm{xor}}(a) \cdot \delta_m^2 \cdot 2^{m-1}
  \qquad \text{as } m \to \infty,
\]
equivalently $\rho_m(a) \to S_{\mathrm{xor}}(a)$.
\end{conjecture}

Conjecture~\ref{conj:main} is the XOR analog of the Hardy--Littlewood
conjecture for arithmetic prime pairs \cite{hardylittlewood1923}:
both predict prime-pair counts as a product over local sieve obstructions.
The arithmetic version predicts $|\{p \leq x : p + 2k \text{ prime}\}|$
via $\prod_q \sigma_q(2k)$ where $\sigma_q$ has a closed residue formula;
our version predicts XOR pair counts via $\prod_q \varphi_q(a)$
where local factors are determined by binary arithmetic.

%% ---------------------------------------------------------------
\section{Numerical Validation}
%% ---------------------------------------------------------------

We compute $C_m(a)$ for all parity-preserving masks of weight $1$--$4$
at $m \in \{16, 18, 20\}$ using the WHT, and compute $S_{\mathrm{xor}}(a)$
by numerical integration of the local factors over odd integers in $[1, 2^{18})$.

\begin{table}[h]
\centering
\caption{Singular series validation: Pearson $r$ between $\rho_m(a)$ and
$S_{\mathrm{xor}}(a)$, and mean ratio $\rho_m/S_{\mathrm{xor}}$,
over all weight $1$--$4$ masks.}
\label{tab:validation}
\begin{tabular}{rrrrrr}
\toprule
$m$ & Masks & Pearson $r$ & Mean $\rho/S_{\mathrm{xor}}$ & Std \\
\midrule
16 & 1940 & 0.886 & 1.001 & 0.04 \\
18 & 3213 & 0.859 & 1.001 & 0.04 \\
20 & 5035 & 0.879 & 1.002 & 0.04 \\
\bottomrule
\end{tabular}
\end{table}

Per-weight breakdown at $m = 20$:

\begin{table}[h]
\centering
\caption{Mean $\rho_m(a)$ vs.\ mean $S_{\mathrm{xor}}(a)$ and ratio by mask weight, $m=20$.}
\label{tab:byweight}
\begin{tabular}{rrrrrr}
\toprule
Weight & $n$ & Mean $\rho$ & Mean $S_{\mathrm{xor}}$ & Mean $\rho/S_{\mathrm{xor}}$ \\
\midrule
1 & 19   & 0.6624 & 0.6660 & 0.9946 \\
2 & 171  & 1.1393 & 1.1398 & 0.9995 \\
3 & 969  & 0.9382 & 0.9347 & 1.0046 \\
4 & 3876 & 1.0358 & 1.0341 & 1.0018 \\
\bottomrule
\end{tabular}
\end{table}

The weight-1 mean $\rho \approx 0.6624$ is consistent with
$S_{\mathrm{xor}}(2^j) = C_2 \approx 0.66016$
(Corollary~\ref{cor:C2}), with the small residual
attributable to finite-$m$ boundary effects.

The local factor formulas of Theorems \ref{thm:weight1} and \ref{thm:weight2}
are verified to six decimal places: the formula $q(q-2)/(q-1)^2$ matches
the numerically computed $\varphi_q(2^j)$ for all primes $q \leq 47$
and all bit positions $j \in \{1, \ldots, 19\}$ with zero spread across $j$.
The weight-2 formula is verified for all primes $q \leq 13$
and all bit-differences $d \in \{1, \ldots, 13\}$.

%% ---------------------------------------------------------------
\section{Discussion}
%% ---------------------------------------------------------------

\paragraph{Scale-dependent structure.}
The XOR autocorrelation reveals a three-scale picture.
At the scale of Hamming balls (radius $\geq 2$),
prime counts are indistinguishable from a random subset:
z-scores follow $N(0,1)$ and low-degree Walsh--Fourier coefficients
are consistent with random baselines.
At the scale of XOR masks of weight $1$--$4$,
a structured sieve fingerprint appears: $\rho_m(a)$ departs from 1
by an amount well-predicted by $S_{\mathrm{xor}}(a)$,
following the alternating pattern
$\rho(k) \approx 1 + C_0(-r)^k$ ($C_0 \approx 0.70$, $r \approx 0.46$)
where $k$ is the mask weight.
At weight $\geq 5$, $|\rho - 1| < 1\sigma$ and primes appear random again.

\paragraph{Comparison with arithmetic Hardy--Littlewood.}
The arithmetic local factor for prime pairs $(p, p+2k)$ at a prime $q \nmid 2k$
is also $q(q-2)/(q-1)^2$ --- identical to our weight-1 formula.
In both cases the factor counts pairs of distinct residues modulo $q$;
the coincidence follows from the shared elementary structure,
not from any deep connection between XOR pairs and arithmetic pairs.

\paragraph{Open problems.}
(i) Does $\rho_m(a)/S_{\mathrm{xor}}(a) \to 1$ exactly, or is there a
persistent $O(1/m)$ correction analogous to the error term in the PNT?
(ii) Do weight-3 local factors admit a closed form in terms of the three
pairwise bit-differences?
(iii) Can Conjecture~\ref{conj:main} be proved for any infinite family of masks?

%% ---------------------------------------------------------------

\begin{thebibliography}{9}

\bibitem{hardylittlewood1923}
G.~H. Hardy and J.~E. Littlewood,
\textit{Some problems of `Partitio Numerorum' III: On the expression of a number as a sum of primes},
Acta Mathematica \textbf{44} (1923), 1--70.

\bibitem{odonnell2014}
R. O'Donnell,
\textit{Analysis of Boolean Functions},
Cambridge University Press, 2014.

\bibitem{granville1995}
A. Granville,
Harald Cram\'er and the distribution of prime numbers,
\textit{Scandinavian Actuarial Journal} \textbf{1995}(1) (1995), 12--28.

\end{thebibliography}

\end{document}
